\xobservationheading\label{sec:observation}

We began revising this paper when we found that the system we had described
was not yet the system we thought we had built. The catalogue named an
active-inference loop, its implementation supplied working components, and
individual episodes supplied evidence of their use. What we had not established
was whether those components, connected as claimed, performed the work that
the description led us to expect. Writing the account made that missing
warrant difficult to ignore.\worknote{IMPROVED: a witness registry now binds paper quantities to Lean
propositions checked at pinned production records --- 35 claims across
eleven node families, green 2026-09-13; the whole-loop machinery closed
grounded on a preregistered cohort (r8, 2026-09-14).
\wkflag{OUTSTANDING}: the whole-loop warrant (rows 24--28) still awaits
the qualifying run (row 28); the certificate reading stands at the
honest-false stage until it runs.}

The immediate trigger was a question about the Learn stage: who actually
consumed its output? The investigation found that an input had stopped arriving,
an observer had stopped running, and writeback deposits had ceased. No standing
monitor had reported the broken handoff. These findings, and the deeper
examination of claims, producers, consumers and their evidence that followed,
are documented in the Futon 2026 working paper, \emph{What Problems Are We
Solving?}. Subsequent runs tested parts of the emerging repair programme;
they were not the sole source of the revision.\worknote{IMPROVED: the stopped loop is startable on demand again --- a
reviewed entrypoint ran end to end 2026-09-13 (run r1); no cron
reinstated. r1's cohort was already at target so no fresh attempt started
(the stopping rule held, correctly); a fresh cohort is staged for r2, and
rows 13/14/15/23 evidence waits on it.}

This matters especially for R20, \emph{Interoceptive Tripwires}
(Section~\ref{pat:r20}). We had already described a mechanism intended to make
failures of the machinery observable. Its presence in the catalogue did not
establish that the necessary observations existed, that a monitor consumed
them, or that its checks could detect a mistaken account of the system. The
audit therefore had to examine the proposed assurance mechanism as well as
the machinery it was meant to observe. Revising R20 raised obligations for
other entries: what they produce, what consumes those products, which identities
and assumptions their handoffs preserve, and what evidence could contradict
their claims.\worknote{IMPROVED: R20 is now a demonstrated live sense: at r1 start
(2026-09-13) wire T10 caught real loaded-code drift, four trips retained.
\wkflag{OUTSTANDING}: the trip-to-commitment-dial link (R20 to R14) --- constants
ruled, integration unbuilt.}

The question of validity is consequently larger than any one entry. Does the
implementation faithfully realise the declared model? Does that model address
the intended problem, rather than a simpler substitute? Evidence for the first
question does not settle the second. In this revision we use the audit to ask
what makes a catalogue answerable to both questions. We do not treat the
existence of tripwires, passing component tests, or the completion of repairs as
validation of the whole system. The catalogue provides the proposed mechanisms;
the investigation shows where their warrants need to be strengthened.\worknote{
Both validity questions now have named instruments: implementation-vs-model
is the per-node witness claim (registry, 35 admitted);
model-vs-problem is the qualifying-run definition with its two-level
certificate reading (ruled 2026-09-13: attest the complete typed census,
plus a separate qualifying predicate that may honestly be false).
\wkflag{OUTSTANDING}: the qualifying run itself, tracker row 28.}
